\documentclass[a4paper,12pt]{article}
\usepackage[utf8]{inputenc}
\usepackage[T1]{fontenc}
\usepackage[italian]{babel}
\usepackage{geometry}
\usepackage{amsmath}
\usepackage{cite}
\usepackage{microtype} % Migliora la giustificazione e previene "overfull hboxes" (testo che esce dai margini)
\usepackage{url}       % Gestisce meglio i link e le citazioni lunghe
\usepackage{hyperref}  % Rende i riferimenti cliccabili

% Impostazioni margini ottimizzate per tesi
\geometry{
 a4paper,
 total={170mm,257mm},
 left=30mm, % Margine sinistro leggermente più ampio per rilegatura
 right=25mm,
 top=25mm,
 bottom=25mm
}

\begin{document}

\section{Stato dell'Arte}

\subsection{Panoramica sugli UAV VTOL}

\subsubsection{Configurazioni principali}
La categoria degli UAV a decollo e atterraggio verticale (VTOL) è vasta e si articola principalmente in tre configurazioni strutturali, ognuna con specifici vantaggi e limitazioni operative:
\begin{itemize}
    \item \textbf{Multirotori:} Questa classe, che include quadricotteri ed esacotteri, è caratterizzata da una notevole semplicità meccanica e dalla capacità di mantenere un \textit{hovering} (volo stazionario) estremamente stabile. Tuttavia, soffrono di una bassa efficienza energetica nel volo traslato, il che ne limita drasticamente l'autonomia e il raggio d'azione.
    \item \textbf{Ala fissa:} I velivoli ad ala fissa convenzionali offrono un'elevata efficienza aerodinamica, permettendo lunghe percorrenze e velocità di crociera elevate. Il loro svantaggio principale risiede nella necessità di infrastrutture dedicate (piste) o spazi ampi per le manovre di decollo e atterraggio, precludendone l'uso in ambienti confinati.
    \item \textbf{Ibridi:} Questa categoria nasce per combinare la versatilità operativa dei multirotori con l'efficienza dei velivoli ad ala fissa. Include architetture complesse come i \textit{tilting-rotor} (rotori basculanti), \textit{tilting-wing} (ali basculanti) e \textit{tailsitter} (velivoli che decollano sulla coda), capaci di transire tra le modalità di volo verticale e orizzontale.
\end{itemize}

\subsubsection{Sfide di modellazione e controllo}
Lo sviluppo di velivoli VTOL ibridi presenta sfide ingegneristiche significative, concentrate prevalentemente nella gestione della fase di transizione tra il volo verticale e quello orizzontale. In questo regime, il velivolo è soggetto a forti non linearità aerodinamiche causate dall'interazione complessa tra i flussi dei rotori, le superfici alari e la fusoliera.
Una modellazione efficace deve catturare accuratamente fenomeni quali gli effetti giroscopici, le variazioni repentine di portanza e resistenza, e il forte accoppiamento tra i gradi di libertà longitudinale e laterale. Di conseguenza, i sistemi di controllo devono essere progettati per garantire robustezza e prestazioni elevate in tutto l'inviluppo di volo, adattandosi a condizioni operative che variano drasticamente tra l'hovering statico e la crociera aerodinamica.

\subsection{UAV VTOL con ali indipendentemente inclinabili}
Un contributo significativo nello sviluppo di configurazioni ibride innovative è stato presentato in [1], dove viene proposto un UAV VTOL caratterizzato da ali inclinabili in modo indipendente. L'obiettivo primario di tale design è superare i limiti di manovrabilità dei tricotteri classici e mitigare la coppia di reazione tipicamente gestita dal rotore di coda.

Dal punto di vista della modellazione, è stato derivato un modello dinamico a sei gradi di libertà (6-DOF) basato sulla formulazione di Newton-Eulero. Il modello evidenzia come il controllo differenziale dell'angolo di inclinazione (\textit{tilt}) delle due ali generi una differenza di portanza in grado di indurre momenti di rollio, sfruttando le forze propulsive in combinazione con quelle aerodinamiche.

La strategia di controllo implementata si basa su regolatori PID e PI integrati nel firmware Ardupilot, strutturati in base alla fase di volo:
\begin{itemize}
    \item \textbf{Modalità Verticale:} Vengono impiegati PID classici per la stabilizzazione di rollio, beccheggio e imbardata, allocando i comandi su tre motori brushless e tre servocomandi.
    \item \textbf{Modalità Orizzontale:} Il controllo del rollio è affidato alla variazione differenziale del tilt alare, il beccheggio è gestito tramite il rotore di coda (spinta e inclinazione vettoriale), mentre l'imbardata è stabilizzata da un controllore PI che minimizza lo scivolamento laterale (\textit{sideslip}).
\end{itemize}
La transizione è governata da una logica a stati finiti che coordina l'inclinazione progressiva delle ali, garantendo una traiettoria fluida durante il cambio di configurazione.

\subsection{Modellazione e controllo Backstepping per Quadricotteri}
In [2], viene affrontato il problema del controllo non lineare per quadricotteri attraverso la derivazione di un modello completo a 6-DOF che include effetti giroscopici e dinamica degli attuatori. L'analisi sottolinea la natura sottattuata e fortemente accoppiata del sistema.
La strategia di controllo proposta adotta un approccio sequenziale ibrido:
\begin{itemize}
    \item Per il \textbf{sottosistema traslazionale} (posizione e altitudine), viene utilizzata una tecnica di linearizzazione con retroazione accoppiata a un regolatore PD.
    \item Per il \textbf{sottosistema rotazionale} (assetto), viene implementato un controllore non lineare basato sulla tecnica \textit{Backstepping} integrata con un'azione integrale (PID).
\end{itemize}
L'uso del Backstepping garantisce la stabilità asintotica globale (dimostrata tramite funzioni di Lyapunov) e offre una robustezza superiore rispetto ai PID tradizionali, specialmente nella reiezione dei disturbi esterni e nella riduzione dell'errore di inseguimento a regime.

\subsection{Disaccoppiamento e Controllo Sliding Mode (SMC)}
Il problema dell'accoppiamento dinamico nei VTOL è trattato specificamente in [3]. Gli autori propongono una procedura di \textit{decoupling} in quattro fasi per trasformare le equazioni del moto, fortemente interconnesse, in una forma standard sottattuata più semplice da controllare.
Su questo modello disaccoppiato viene applicato un controllore \textit{Sliding Mode Control} (SMC). La scelta dello SMC è motivata dalla sua intrinseca capacità di gestire incertezze parametriche e disturbi esterni mantenendo il sistema sulla superficie di scorrimento desiderata. Le simulazioni confermano che l'integrazione tra la tecnica di disaccoppiamento e il controllo robusto SMC permette di annullare l'errore di tracciamento in tempo finito, eliminando al contempo il fenomeno del \textit{chattering} (vibrazioni ad alta frequenza) tipico di queste strategie.

\subsection{Modellazione e Controllo Robusto per UAV Coaxial Tilt-Rotor}
Un'evoluzione significativa nel panorama dei velivoli a decollo verticale è rappresentata dalla configurazione \textit{Coaxial Tilt-Rotor} (CTRUAV), analizzata approfonditamente in [4]. Questa architettura si distingue per la presenza di due coppie di rotori coassiali inclinabili posizionati nella parte anteriore del velivolo e un rotore fisso in coda. Rispetto ai tricotteri convenzionali, l'impiego di rotori coassiali permette di generare una spinta maggiore senza aumentare le dimensioni del telaio, migliorando la manovrabilità e annullando naturalmente il momento giroscopico grazie alla configurazione controrotante.

Il modello dinamico a sei gradi di libertà è derivato utilizzando le formule di Newton-Eulero. Per governare questo sistema in presenza di disturbi esterni, è stata proposta una strategia di controllo gerarchica:
\begin{itemize}
    \item \textbf{Loop Esterno (Posizione):} Per il tracciamento della traiettoria desiderata, viene impiegato un controllore basato sulla tecnica del \textit{Backstepping}, che stabilizza asintoticamente l'errore di posizione.
    \item \textbf{Loop Interno (Assetto):} La stabilizzazione è affidata a un controllore \textit{Robust Adaptive Backstepping Sliding Mode} (ABSMC). L'integrazione di una legge adattiva permette di stimare online il limite superiore dei disturbi esterni. Tale stima modula il guadagno della legge di controllo, minimizzando il fenomeno del \textit{chattering} tramite l'uso di funzioni di saturazione.
\end{itemize}
Le simulazioni dimostrano che l'approccio proposto garantisce una robustezza superiore rispetto ai controllori backstepping tradizionali.

\subsection{Strategie di Volo e Allocazione del Controllo per Tricotteri Tilt-Rotor}
In [5], l'attenzione si sposta sulla realizzazione pratica di un volo autonomo "full-mode" per un UAV tilt-rotor in configurazione tricottero compatta. L'obiettivo è superare le limitazioni di costo associate all'ottenimento di modelli dinamici precisi, proponendo una strategia robusta che non dipenda da dati aerodinamici ad alta fedeltà.

La gestione del volo è affidata a una strategia basata su una macchina a stati finiti che coordina le diverse fasi: \textit{Rotor mode} (decollo verticale), \textit{Conversion} (transizione accelerata), \textit{Fixed-wing mode} (crociera) e \textit{Reconversion} (ritorno all'hovering).
Un elemento centrale è il sistema di \textbf{Control Allocation (Mixer)}. Poiché il sistema è sovra-attuato e non affine durante la transizione, il mixer distribuisce i pesi di controllo tra modalità rotore e ala fissa in base all'efficienza degli attuatori e all'angolo di tilt corrente. Una peculiarità è il controllo dell'imbardata in hovering tramite tilt differenziale dei rotori anteriori, che viene disattivato progressivamente durante la transizione.
[cite_start]I test di volo reali hanno validato la capacità del velivolo di eseguire l'intero inviluppo di volo mantenendo la stabilità[cite: 5].

\subsection{Controllo Sliding Mode a Tempo Fisso (Fixed-Time)}
La problematica della convergenza rapida e garantita è affrontata in [6], dove viene proposto uno schema per la stabilizzazione dell'assetto di un \textit{Tilt Trirotor UAV} basato sulla stabilità a tempo fisso (\textit{Fixed-Time Stability}). A differenza della stabilità asintotica o a tempo finito, questo approccio garantisce che il sistema raggiunga l'equilibrio entro un tempo limitato, noto a priori e indipendente dallo stato iniziale.

Il contributo teorico principale è la formulazione di un \textit{Novel Fixed-Time Stable System} che dimostra un tasso di convergenza superiore rispetto ai metodi esistenti. Basandosi su questo, è stato sviluppato uno schema \textit{Nonsingular Fixed-Time Sliding Mode Control} (NFTSMC) che include:
\begin{itemize}
    \item \textbf{Superficie di Scorrimento Non Singolare:} Progettata per evitare le singolarità matematiche tipiche dei controllori terminali e garantire la convergenza.
    \item \textbf{Tecnica Adattiva per i Disturbi:} Una legge di aggiornamento stima i disturbi in tempo reale, assicurando robustezza.
\end{itemize}
[cite_start]La validazione sperimentale su prototipo reale ha confermato una convergenza degli errori di assetto in meno di 0,3 secondi e una notevole reiezione dei disturbi[cite: 6].

\subsection{Allocazione del Controllo tramite Spinta Differenziale}
Un aspetto tecnologico cruciale esplorato in [7] è la completa sostituzione delle superfici di controllo aerodinamiche convenzionali (alettoni, elevatori, timoni) con una logica di controllo basata esclusivamente sulla modulazione della spinta (\textit{Differential Thrust}).
Questa architettura richiede l'impiego di propulsori in grado di fornire spinta reversibile (positiva e negativa), permettendo la generazione di momenti di controllo anche a velocità nulla, una capacità preclusa ai sistemi aerodinamici tradizionali.

La strategia di allocazione del controllo (\textit{Control Allocation}) mappa i comandi di assetto desiderati (Rollio, Beccheggio, Imbardata) direttamente sulle coppie di motori, secondo lo schema seguente (riferito a una configurazione a quattro propulsori disposti a "X"):

\begin{itemize}
    \item \textbf{Controllo di Beccheggio (Pitch):} Sostituisce la funzione dell'elevatore. Viene generato creando una spinta differenziale tra i rotori superiori e quelli inferiori. Ad esempio, per cabrare (muso in su), i propulsori inferiori generano spinta positiva (avanti) mentre quelli superiori generano spinta negativa (indietro), creando una coppia pura sull'asse trasversale.
    
    \item \textbf{Controllo di Imbardata (Yaw):} Sostituisce il timone verticale. Si ottiene differenziando la spinta tra i motori di destra e quelli di sinistra. Una rotazione verso destra è indotta da una spinta positiva sui motori di sinistra e negativa su quelli di destra.
    
    \item \textbf{Controllo di Rollio (Roll):} Sostituisce gli alettoni. A differenza dei velivoli ad ala fissa dove il rollio è generato alle estremità alari, qui viene prodotto tramite una configurazione incrociata (diagonale). [cite_start]Ad esempio, una spinta positiva del motore in basso a destra combinata con il motore in alto a sinistra, opposta alla coppia degli altri due, genera il momento torcente longitudinale necessario[cite: 7].
\end{itemize}

Questa metodologia offre il vantaggio di eliminare attuatori meccanici complessi (servocomandi) e parti mobili, riducendo il peso strutturale e la manutenzione. Tuttavia, introduce una maggiore complessità nel software di controllo, poiché i gradi di libertà risultano fortemente accoppiati: una singola manovra richiede spesso la modulazione simultanea di tutti i propulsori per evitare effetti secondari indesiderati sugli altri assi.

\begin{thebibliography}{99}

\bibitem{1}
K.-J. Nam, J. Joung, and D. Har, ``Tri-copter UAV with individually tilted main wings for flight maneuvers,'' \textit{IEEE Access}, vol. 8, pp. 46753–46770, 2020.

\bibitem{2}
A. A. Mian and D. Wang, ``Modeling and Backstepping-based Nonlinear Control Strategy for a 6 DOF Quadrotor Helicopter,'' \textit{Chinese Journal of Aeronautics}, vol. 21, no. 3, pp. 261–268, 2008.

\bibitem{3}
Y. Zhang, F. Gao, and T. Y. Zeng, ``A sliding mode controller design for vertical take-off and landing based on model decoupling,'' in \textit{Proc. 33rd Chinese Control Conference}, pp. 115–119, 2014.

\bibitem{4}
W.-G. Cheng, R. Wang, and Y.-H. Li, ``Robust backstepping sliding mode control for coaxial tilt-rotor UAV,'' in \textit{2022 5th International Symposium on Autonomous Systems (ISAS)}, IEEE, 2022.

\bibitem{5}
C. Chen, J. Zhang, D. Zhang, and L. Shen, ``Control and flight test of a tilt-rotor unmanned aerial vehicle,'' \textit{International Journal of Advanced Robotic Systems}, vol. 14, no. 1, 2017.

\bibitem{6}
L. Yu, G. He, X. Wang, and S. Zhao, ``Robust Fixed-Time Sliding Mode Attitude Control of Tilt Trirotor UAV in Helicopter Mode,'' \textit{IEEE Transactions on Industrial Electronics}, vol. 69, no. 10, pp. 10322–10332, Oct. 2022.

\bibitem{7}
C. E. D. Riboldi and A. Rolando, ``Thrust-Based Flight Stabilization and Guidance for Autonomous Airships,'' \textit{Aerospace Europe Conference 2023 - 10th EUCASS - 9th CEAS}, 2023. DOI: 10.13009/EUCASS2023-025.

\end{thebibliography}

\end{document}